\begin{solapartat}
\punts{0,5} Desintegració $\beta^-$ \quad ${}^{40}_{19}K \to {}^{40}_{20}Ca + \beta^- + {}^{0}_{0}\bar{\nu}_e$\\
Si no es posa l'antineutrí: \punts{-0,2}
% A l'original aquesta penalització s'escriu «-0-2p» (errata per «-0,2 p», com la del neutrí).

\punts{0,5} Desintegració $\beta^+$ \quad ${}^{40}_{19}K \to {}^{40}_{18}Ar + \beta^+ + {}^{0}_{0}\nu_e$\\
Si no es posa el neutrí: \punts{-0,2}

\[ \beta^- = {}^{0}_{-1}e\ ;\ \beta^+ = {}^{0}_{1}e \]
\end{solapartat}

\begin{solapartat}
% DUBTE: l'enunciat i la pauta diuen 1460 MeV; l'energia real del raig gamma del 40K és 1460 keV (1,460 MeV). Es transcriu tal com és.
\punts{0,1} $E = 1460\ \mathrm{MeV} = 1460\cdot 10^{6}\ \mathrm{eV}\cdot\dfrac{1,6\cdot 10^{-19}}{1\ \mathrm{eV}} = 2,336\cdot 10^{-10}\ \mathrm{J}$

\punts{0,2} $E = h\nu \Rightarrow \nu = \dfrac{E}{h} = \dfrac{2,336\cdot 10^{-10}}{6,63\cdot 10^{-34}} = 3,52\cdot 10^{23}\ \mathrm{s^{-1}}$

\punts{0,3} $\lambda = \dfrac{c}{\nu} = \dfrac{3\cdot 10^{8}}{3,52\cdot 10^{23}} = 8,51\cdot 10^{-16}\ \mathrm{m}$

\punts{0,4} $\Delta m = \dfrac{\Delta E}{c^2} = \dfrac{2,336\cdot 10^{-10}}{(3,00\cdot 10^{8})^2} = 2,60\cdot 10^{-27}\ \mathrm{kg}$
\end{solapartat}
